\documentclass[10pt,a4paper]{amsart}
\usepackage{amsmath,amssymb,amsthm,graphicx}
\usepackage{verbatim}
\newenvironment{reference}{\comment}{\endcomment}
%\newenvironment{reference}{}{}
\newenvironment{slogan}{\comment}{\endcomment}
\newenvironment{history}{\comment}{\endcomment}

% For commutative diagrams we use Xy-pic
\usepackage[all]{xy}

% We use 2cell for 2-commutative diagrams.
\xyoption{2cell}
% Theorem environments.
%
\theoremstyle{plain}
\newtheorem{theorem}[subsection]{Theorem}
\newtheorem{proposition}[subsection]{Proposition}
\newtheorem{lemma}[subsection]{Lemma}

\theoremstyle{definition}
\newtheorem{definition}[subsection]{Definition}
\newtheorem{example}[subsection]{Example}
\newtheorem{exercise}[subsection]{Exercise}
\newtheorem{situation}[subsection]{Situation}

\theoremstyle{remark}
\newtheorem{remark}[subsection]{Remark}
\newtheorem{remarks}[subsection]{Remarks}

\numberwithin{equation}{subsection}

% Macros
%
\def\lim{\mathop{\mathrm{lim}}\nolimits}
\def\colim{\mathop{\mathrm{colim}}\nolimits}
\def\Spec{\mathop{\mathrm{Spec}}}
\def\Hom{\mathop{\mathrm{Hom}}\nolimits}
\def\Ext{\mathop{\mathrm{Ext}}\nolimits}
\def\SheafHom{\mathop{\mathcal{H}\!\mathit{om}}\nolimits}
\def\SheafExt{\mathop{\mathcal{E}\!\mathit{xt}}\nolimits}
\def\Sch{\mathit{Sch}}
\def\Mor{\mathop{\mathrm{Mor}}\nolimits}
\def\Ob{\mathop{\mathrm{Ob}}\nolimits}
\def\Sh{\mathop{\mathit{Sh}}\nolimits}
\def\NL{\mathop{N\!L}\nolimits}
\def\CH{\mathop{\mathrm{CH}}\nolimits}
\def\proetale{{pro\text{-}\acute{e}tale}}
\def\etale{{\acute{e}tale}}
\def\QCoh{\mathit{QCoh}}
\def\Ker{\mathop{\mathrm{Ker}}}
\def\Im{\mathop{\mathrm{Im}}}
\def\Coker{\mathop{\mathrm{Coker}}}
\def\Coim{\mathop{\mathrm{Coim}}}

% Boxtimes
%
\DeclareMathSymbol{\boxtimes}{\mathbin}{AMSa}{"02}

%
% Macros for moduli stacks/spaces
%
\def\QCohstack{\mathcal{QC}\!\mathit{oh}}
\def\Cohstack{\mathcal{C}\!\mathit{oh}}
\def\Spacesstack{\mathcal{S}\!\mathit{paces}}
\def\Quotfunctor{\mathrm{Quot}}
\def\Hilbfunctor{\mathrm{Hilb}}
\def\Curvesstack{\mathcal{C}\!\mathit{urves}}
\def\Polarizedstack{\mathcal{P}\!\mathit{olarized}}
\def\Complexesstack{\mathcal{C}\!\mathit{omplexes}}
% \Pic is the operator that assigns to X its picard group, usage \Pic(X)
% \Picardstack_{X/B} denotes the Picard stack of X over B
% \Picardfunctor_{X/B} denotes the Picard functor of X over B
\def\Pic{\mathop{\mathrm{Pic}}\nolimits}
\def\Picardstack{\mathcal{P}\!\mathit{ic}}
\def\Picardfunctor{\mathrm{Pic}}
\def\Deformationcategory{\mathcal{D}\!\mathit{ef}}
\usepackage{cleveref}
\crefname{originallemma}{original lemma}{original lemmas}
\Crefname{originallemma}{Original lemma}{Original lemmas}
\makeatletter
\newlabel{chow-lemma-chern-classes-E-tensor-L}{{02UD}{source}{}{sourcepointer-1}{}}
\newlabel{chow-lemma-chern-classes-E-tensor-L@cref}{{[originallemma][0][]02UD}{[1][0][]source}{}{}{}}
\makeatother
\crefname{originallemma}{original lemma}{original lemmas}
\Crefname{originallemma}{Original lemma}{Original lemmas}
\makeatletter
\newlabel{chow-lemma-pushforward-cap-cj}{{02U9}{source}{}{sourcepointer-2}{}}
\newlabel{chow-lemma-pushforward-cap-cj@cref}{{[originallemma][0][]02U9}{[1][0][]source}{}{}{}}
\makeatother
\crefname{originallemma}{original lemma}{original lemmas}
\Crefname{originallemma}{Original lemma}{Original lemmas}
\makeatletter
\newlabel{chow-lemma-top-chern-class}{{0FA9}{source}{}{sourcepointer-3}{}}
\newlabel{chow-lemma-top-chern-class@cref}{{[originallemma][0][]0FA9}{[1][0][]source}{}{}{}}
\makeatother
\crefname{originalreference}{original reference}{original references}
\Crefname{originalreference}{Original reference}{Original references}
\makeatletter
\newlabel{lemma-projective-space-bundle-formula}{{0FGQ}{source}{}{sourcepointer-4}{}}
\newlabel{lemma-projective-space-bundle-formula@cref}{{[originalreference][0][]0FGQ}{[1][0][]source}{}{}{}}
\makeatother
\title{The class of the diagonal of a projective bundle}
\author{The Stacks Project Authors}
\date{stacks-project:weil.tex:local-snapshot-20261012; GNU FDL 1.2 or later}
\begin{document}
\maketitle
\noindent\textit{Editorial note.} All delivered mathematical text is the author's own, copied from the pinned Stacks Project chapter file weil.tex; no statement or proof step was rewritten, reordered or removed.\par\medskip
\noindent\textit{Editorial note.} The counted proof is the author's complete argument: the reduction to the fibre product over the base, the description of the diagonal as the zero scheme of the displayed map, and the Chern class computation for the twist.\par\medskip
\noindent\textit{Editorial note.} The label lemma-diagonal-projective-bundle names the counted statement, and the proof environment immediately following it in the source is delivered verbatim as its proof.\par\medskip
\noindent\textit{Editorial note.} The counted statement is carried by the author's own lemma environment in the chapter source.\par\medskip
\noindent\textit{Editorial note.} Pre-open gate: python operations/agent-authoring/tools/precard\_gate.py --source weil.tex --label lemma-diagonal-projective-bundle --start 1465 --end 1517 --local-id 28143 returned RESULT=OK (L28143 free, occupancy FREE, public interval clear). Pre-assembly tag self-check: canonical tag resolved and not present in reference/public-items.jsonl.\par\medskip
\noindent\textit{Editorial note.} The statement and proof are copied from the pinned chapter file; the cited inputs are kept as external pointers with their canonical Stacks tags.\par\medskip
\noindent\textit{Allowed standard prerequisites (editorial).}\begin{itemize}
\item The class of a subscheme cut out by a regular section as a top Chern class. Chow Homology, the lemma identifying the class of the zero scheme of a regular section with the top Chern class of the dual of the defining module. Expresses the class of the diagonal as the top Chern class of an explicitly described module.
\item Chern classes of a tensor product with an invertible module. Chow Homology, the lemma computing the Chern classes of the tensor product of a locally free module with an invertible module. Turns the top Chern class of the twisted module into the displayed sum over binomial coefficients.
\item Compatibility of capping with Chern classes with proper pushforward. Chow Homology, the lemma stating that capping with Chern classes commutes with proper pushforward. Reduces the statement on the product to the statement on the fibre product over the base.
\end{itemize}\medskip
\ifdefined\hypertarget\hypertarget{sourcepointer-1}{}\fi
\noindent\textit{Source reference.} Original lemma 02UD: sources/stacks-chow/tags.txt, source lines 3728--3728. Prerequisite pointer invoked at proof line(s) 1516; it names a standard fact of a sibling chapter that the counted argument uses, and every step of the counted proof built on it is retained verbatim in the delivered body.\par\medskip
\ifdefined\hypertarget\hypertarget{sourcepointer-2}{}\fi
\noindent\textit{Source reference.} Original lemma 02U9: sources/stacks-chow/tags.txt, source lines 3724--3724. Prerequisite pointer invoked at proof line(s) 1487; it names a standard fact of a sibling chapter that the counted argument uses, and every step of the counted proof built on it is retained verbatim in the delivered body.\par\medskip
\ifdefined\hypertarget\hypertarget{sourcepointer-3}{}\fi
\noindent\textit{Source reference.} Original lemma 0FA9: sources/stacks-chow/tags.txt, source lines 19104--19104. Prerequisite pointer invoked at proof line(s) 1510; it names a standard fact of a sibling chapter that the counted argument uses, and every step of the counted proof built on it is retained verbatim in the delivered body.\par\medskip
\ifdefined\hypertarget\hypertarget{sourcepointer-4}{}\fi
\noindent\textit{Source reference.} Original reference 0FGQ: sources/stacks-weil/pinned-original-source.tex, source lines 1391--1391. Prerequisite pointer invoked at proof line(s) 1467; it names a standard fact of the same chapter that the counted argument uses, and every step of the counted proof built on it is retained verbatim in the delivered body.\par\medskip
\begin{lemma}
\label{lemma-diagonal-projective-bundle}
Let $p : P \to X$ be as in Lemma \ref{lemma-projective-space-bundle-formula}.
The class $[\Delta_P]$ of the diagonal of $P$ in $\CH^*(P \times P)$
can be written
as
$$
[\Delta_P] =
\left(\sum\nolimits_{i = 0, \ldots, r - 1}
{r - 1 \choose i} c_{r - 1 - i}(\text{pr}_1^*\mathcal{S}^\vee) \cap
c_1(\text{pr}_2^*\mathcal{O}_P(1))^i\right)
\cap
(p \times p)^*[\Delta_X]
$$
where $\mathcal{S}$ is the kernel of the canonical surjection
$p^*\mathcal{E} \to \mathcal{O}_P(1)$.
\end{lemma}

\begin{proof}
Observe that $(p \times p)^*[\Delta_X] = [P \times_X P]$.
Since $\Delta_P \subset P \times_X P \subset P \times P$
and since capping with Chern classes commutes with proper pushforward
(Chow Homology, Lemma \ref{chow-lemma-pushforward-cap-cj})
it suffices to show that the class of
$\Delta_P \subset P \times_X P$ in $\CH^*(P \times_X P)$
is equal to
$$
\left(\sum\nolimits_{i = 0, \ldots, r - 1}
{r - 1 \choose i} c_{r - 1 - i}(q_1^*\mathcal{S}^\vee) \cap
c_1(q_2^*\mathcal{O}_P(1))^i\right)
\cap
[P \times_X P]
$$
where $q_i : P \times_X P \to P$, $i = 1, 2$ are the projections.
Set $q = p \circ q_1 = p \circ q_2 : P \times_X P \to X$.
Consider the maps
$$
q_1^*\mathcal{S} \otimes q_2^*\mathcal{O}_P(-1) \to
q^*\mathcal{E} \otimes q^*\mathcal{E}^\vee \to
\mathcal{O}_{P \times_X P}
$$
where the final arrow is the pullback by $q$ of the evaluation map
$\mathcal{E} \otimes_{\mathcal{O}_X} \mathcal{E}^\vee \to \mathcal{O}_X$.
The source of the composition is a module locally free of rank $r - 1$
and a local calculation shows that this map vanishes exactly along
$\Delta_P$. By Chow Homology, Lemma \ref{chow-lemma-top-chern-class}
the class $[\Delta_P]$ is the top Chern class of the dual
$$
q_1^*\mathcal{S}^\vee \otimes q_2^*\mathcal{O}_P(1)
$$
The desired result follows from Chow Homology, Lemma
\ref{chow-lemma-chern-classes-E-tensor-L}.
\end{proof}

\end{document}
